\chapter{Symbolic validation of \texorpdfstring{$h_{22}$}{h22} and
         \texorpdfstring{$h_{21}$}{h21}}
\label{ap:symb-val}

This appendix presents the symbolic validation of the \textsc{ARES}
expansion coefficients $h_{22}$ and $h_{21}$ against the corresponding
SCET subtraction coefficients $\mathcal{C}^{(2)}_1$ and
$\mathcal{C}^{(2)}_0$. The comparison is performed using
\textsc{Mathematica}, and the complete notebook is available in
Ref.~\cite{peake2026validation}.

The notation used in the \textsc{Mathematica} implementation is summarised
in Table~\ref{tab:mma-symbols}. Uppercase symbols denote the full
expansion of a subtraction ingredient, following
Eq.~\eqref{eq:subt-mn-exp}. The lowercase symbols are the corresponding
coefficients defined in Appendix~\ref{ap:ARES-coeff}.

\begin{table}[h]
\begin{center}
\begin{tabular}{ll|ll}
\hline\hline
Notebook & LaTeX & Notebook & LaTeX \\
\hline\hline
\texttt{J[m,n]} & $J^{(m)}_{\kappa_3,n}$ &
\texttt{Ck[l]} & $C_\ell$ \\
\texttt{B1[m,n]}, \texttt{B2[m,n]} &
$B^{(m)}_{\kappa_{1,2},n}/f_{i,j}$ &
\texttt{Ctot} & $\sum_\ell C_\ell$ \\
\texttt{S[m,n]} & $\widehat{S}^{(m)}_{\kappa,n}$ &
\texttt{g[l,m]} & $\gamma^{(m)}_\ell$ (SCET) \\
\texttt{Q[l]}, \texttt{Q[l,lbar]} &
$Q_\ell$, $Q_{\ell\bar\ell}$ &
\texttt{gA[l,m]} & $\gamma^{(m)}_\ell$ (ARES) \\
\texttt{Qresum} & $Q$ &
\texttt{gs1} & $\gamma^{(1)}_S$ \\
\texttt{H1SCET} & $H^{(1)}_{\rm SCET}$ &
\texttt{P[l,m]} & $P^{(m)}_{ik}(z)$ \\
\hline\hline
\end{tabular}
\end{center}
\caption{Correspondence between the symbols used in the
\textsc{Mathematica} notebook and the notation of this thesis.}
\label{tab:mma-symbols}
\end{table}

For compactness, we define the following shorthand:
\begin{lstlisting}[style=mma, caption={Definition of the shorthand used
for the subtraction ingredients.}]
tot[m_, n_] := J[m, n] + B1[m, n] + B2[m, n] + S[m, n];
\end{lstlisting}

The subtraction coefficients $\mathcal{C}^{(2)}_1$ and
$\mathcal{C}^{(2)}_0$ are implemented directly from
Eqs.~\eqref{eq:C21} and~\eqref{eq:C20}:
\begin{lstlisting}[style=mma, caption={Implementation of
$\mathcal{C}^{(2)}_1$ and $\mathcal{C}^{(2)}_0$.}]
C21 = (tot[2, 1] + H1SCET tot[1, 1]
    + 2 (B1[1, 0] J[1, 0] + B1[1, 0] B2[1, 0] + ... )
    - \[Pi]^2/3 (B1[1, 1] J[1, 1] + B1[1, 1] B2[1, 1] + ... )
    + B1[1, -1] J[1, 1] + B2[1, -1] J[1, 1] + ... );

C20 = (tot[2, 0] + H1SCET tot[1, 0]
    - \[Pi]^2/6 (B1[1, 0] J[1, 1] + B2[1, 0] J[1, 1] + ... )
    + 2 Zeta[3] (B1[1, 1] J[1, 1] + B1[1, 1] B2[1, 1] + ... )
    + B1[1, -1] J[1, 0] + B2[1, -1] J[1, 0] + ... );
\end{lstlisting}

The \textsc{ARES} coefficients are defined according to
Eqs.~\eqref{eq:h11}, \eqref{eq:h22}, and~\eqref{eq:h21}. Since the
cross-terms $h_{10}h_{12}$ and $h_{10}h_{11}$ have already been validated
at lower logarithmic accuracy, they are omitted here. The same terms are
removed from the corresponding SCET coefficients to simplify the
comparison:
\begin{lstlisting}[style=mma, caption={Removal of the cross-terms already
validated in the \textsc{ARES} coefficients.}]
C11  = tot[1, 1];
C10  = tot[1, 0];
C1m1 = tot[1, -1] + H1SCET;

C21Rem = C21 - C11 C1m1;
C20Rem = C20 - C10 C1m1;
\end{lstlisting}

The resulting \textsc{ARES} expressions are implemented as
\begin{lstlisting}[style=mma, caption={Implementation of the \textsc{ARES}
expansion coefficients $h_{11}$, $h_{22}$ and $h_{21}$.}]
h11 = ((g[1, 0] + g[2, 0] + g[3, 0])
    - 2 Ck[1] Log[(Q[1, 2]^2 Q[1, 3]^2)/(Q[2, 3]^2 M Q[1])] + ...
    - P[1, 0] - P[2, 0]);

h22 = (h11^2/2 - \[Pi]^2/3 Ctot^2 - K1 Ctot
    + 2 \[Pi] \[Beta]0 (-2 (Ck[1] + Ck[2] - Ck[3])
        Log[Q[1, 2]^2/Qresum^2]
        + g[1, 0]/2 - 3 Ck[1] Log[Qresum^2/(M Q[1])]
        - 1/2 P[1, 0] + ... )
    + 4 \[Pi] Ctot \[Beta]0 Log[Qresum/\[Mu]R]);

h21 = (-K1 ((Ck[1] + Ck[2] - Ck[3])
             (Log[Q[1, 2]^2/(M Q[1])]
              + Log[Q[1, 2]^2/(M Q[2])) + ... )
    + 2 \[Pi] \[Beta]0 ( ... )
    + 2 Ctot^2 PolyGamma[2, 1]
    - gA[1, 1] - gA[2, 1] - gA[3, 1]
    - (P[1, 1] + P[2, 1]));
\end{lstlisting}

Using Eq.~\eqref{eq:coeff-convent}, $\mathcal{C}^{(2)}_1=8h_{22}$.
The symbolic comparison gives
\begin{lstlisting}[style=mma, caption={Symbolic comparison of $h_{22}$ with
$\mathcal{C}^{(2)}_1$.}]
C21Rem - 8 h22 // FullSimplify

(* Out: 0 *)
\end{lstlisting}
showing exact agreement.

For $h_{21}$, Eq.~\eqref{eq:coeff-convent} gives
$\mathcal{C}^{(2)}_0=-4h_{21}$. The symbolic comparison gives
\begin{lstlisting}[style=mma, caption={Symbolic comparison of $h_{21}$ with
$\mathcal{C}^{(2)}_0$.}]
g1Res = C20Rem + 4 h21 // FullSimplify

(* Out: (1/2)(-g[1, 1] - 8 gA[1, 1] - g[2, 1] - 8 gA[2, 1]
               - g[3, 1] - 8 gA[3, 1] - 2 Ctot gs1) *)
\end{lstlisting}

The residual arises from the different conventions for the two-loop
anomalous dimensions in the SCET and \textsc{ARES} expressions, as
discussed in Eqs.~\eqref{eq:gammal1-flav},
\eqref{eq:gamma-1-SCET-q}, and~\eqref{eq:gamma-1-SCET-g}. Substituting
the explicit values for the $q\bar q\to Zg$ channel gives
\begin{lstlisting}[style=mma, caption={Evaluation of the residual for the
$q\bar q\to Zg$ channel.}]
g1Res /. {g[1, 1] -> ..., gA[1, 1] -> ..., gs1 -> ...,
          Ck[1] -> CF, Ck[2] -> CF, Ck[3] -> CA} // FullSimplify

(* Out: 0 *)
\end{lstlisting}

The remaining partonic channels are validated in the same way, using the
corresponding colour factors and anomalous dimensions.